% ============================================================
% common/questions_env.tex —— 题目 / 解析的语义包装环境
%
% questions.tex 中每道题包在 \begin{question}...\end{question} 里，
% answers.tex 中每题解析包在 \begin{answer}...\end{answer} 里；
% 题头排版仍由内容行中的 \gkti / \gksolhead 完成（见 macros.tex）。
%
% 环境本身不产生任何盒子或垂直间距（只引入一个 TeX 分组），
% 排版结果与不包裹时逐点一致。它的价值在于：
%   1. 给每道题一个机器可数的结构边界（校验脚本数环境数 = 题量）；
%   2. 为将来的条件编译留钩子——某版本可重定义环境实现按题取舍。
% ============================================================

% 环境不把整题装盒：题目与解析仍可自然跨页，但开始题块前必须为题面、
% 图形/表格或解析首段留出足够空间。仅保留 2 个 baseline 会让题干、
% 选项和图形在页边界处被拆开，形成不可打印的孤立内容；题目保留
% 0.15\textheight，在避免短题制造大段空白与避免题块拆分之间取平衡。
% 解析题块保留 0.5\textheight，以容纳“题面 + 图形/选项 + 首段解析”；
% 结构标题仍用更大的保留量；末页自然收尾由视觉复审判定。
\newenvironment{question}{\Needspace{0.15\textheight}}{}
\newenvironment{answer}{\Needspace{0.5\textheight}}{}
